\textbf{Phần 1. Trắc nghiệm lựa chọn}
\setcounterpageref{ex}{0}

\begin{ex} %[5H8V46]
Đổi số đo $4\text{ dm}^3\text{ }25\text{ cm}^3$ ra đề-xi-mét khối ta được số thập phân là
\choice
{$4{,}25\text{ dm}^3$}
{\True $4{,}025\text{ dm}^3$}
{$4{,}0025\text{ dm}^3$}
{$42{,}5\text{ dm}^3$}
\loigiai{
Vì $1\text{ dm}^3 = 1\,000\text{ cm}^3$ nên $25\text{ cm}^3 = 0{,}025\text{ dm}^3$.\\
Do đó: $4\text{ dm}^3\text{ }25\text{ cm}^3 = 4{,}025\text{ dm}^3$.\\
\textbf{Đáp án của câu hỏi là B}
}
\end{ex}

\begin{ex} %[5H8V46]
Một hộp sữa tươi có ghi thể tích thực là $180\text{ ml}$ (biết $1\text{ ml} = 1\text{ cm}^3$). Thể tích của hộp sữa tươi đó tính theo đề-xi-mét khối là
\choice
{$1{,}8\text{ dm}^3$}
{\True $0{,}18\text{ dm}^3$}
{$0{,}018\text{ dm}^3$}
{$18\text{ dm}^3$}
\loigiai{
Ta có: $180\text{ ml} = 180\text{ cm}^3$.\\
Đổi ra đề-xi-mét khối: $180 : 1\,000 = 0{,}18\text{ (dm}^3\text{)}$.\\
\textbf{Đáp án của câu hỏi là B}
}
\end{ex}

\begin{ex} %[5H8V46]
Điền dấu thích hợp vào chỗ chấm: $5\text{ dm}^3\text{ }60\text{ cm}^3 \dots 5{,}6\text{ dm}^3$.
\choice
{\True $<$}
{$>$}
{$=$}
{Không so sánh được}
\loigiai{
Ta đổi: $5\text{ dm}^3\text{ }60\text{ cm}^3 = 5{,}06\text{ dm}^3$.\\
Vì $5{,}06 < 5{,}60$ nên $5\text{ dm}^3\text{ }60\text{ cm}^3 < 5{,}6\text{ dm}^3$.\\
\textbf{Đáp án của câu hỏi là A}
}
\end{ex}

\begin{ex} %[5H8V46]
Một bình nước chứa $2{,}5\text{ l}$ nước. Người ta rót đều lượng nước đó vào $10$ chiếc cốc nhỏ. Mỗi chiếc cốc chứa số xăng-ti-mét khối nước là
\choice
{$25\text{ cm}^3$}
{\True $250\text{ cm}^3$}
{$2\,500\text{ cm}^3$}
{$0{,}25\text{ cm}^3$}
\loigiai{
Đổi: $2{,}5\text{ l} = 2{,}5\text{ dm}^3 = 2\,500\text{ cm}^3$.\\
Mỗi cốc chứa số xăng-ti-mét khối nước là:
$2\,500 : 10 = 250\text{ (cm}^3\text{)}$.\\
\textbf{Đáp án của câu hỏi là B}
}
\end{ex}

\begin{ex} %[5H8V46]
Giá trị của biểu thức $3{,}2\text{ dm}^3 + 450\text{ cm}^3 \times 4$ tính theo đơn vị đề-xi-mét khối là
\choice
{$4{,}5\text{ dm}^3$}
{\True $5\text{ dm}^3$}
{$5{,}5\text{ dm}^3$}
{$50\text{ dm}^3$}
\loigiai{
Ta có: $450\text{ cm}^3 \times 4 = 1\,800\text{ cm}^3 = 1{,}8\text{ dm}^3$.\\
Biểu thức bằng: $3{,}2 + 1{,}8 = 5\text{ (dm}^3\text{)}$.\\
\textbf{Đáp án của câu hỏi là B}
}
\end{ex}

\begin{ex} %[5H8V46]
Một bể cá hình hộp chữ nhật có lòng trong dài $4\text{ dm}$, rộng $2{,}5\text{ dm}$ và cao $3\text{ dm}$. Khi bể chứa đầy nước thì thể tích nước trong bể là bao nhiêu lít?
\choice
{$25\text{ l}$}
{\True $30\text{ l}$}
{$35\text{ l}$}
{$40\text{ l}$}
\loigiai{
Thể tích của bể cá là:
$4 \times 2{,}5 \times 3 = 30\text{ (dm}^3\text{)}$.\\
Vì $1\text{ dm}^3 = 1\text{ l}$ nên bể chứa được $30\text{ l}$ nước.\\
\textbf{Đáp án của câu hỏi là B}
}
\end{ex}

\begin{ex} %[5H8V46]
Một khối kim loại có thể tích $1{,}2\text{ dm}^3$. Biết mỗi xăng-ti-mét khối kim loại đó cân nặng $7{,}8\text{ g}$. Khối kim loại đó cân nặng bao nhiêu ki-lô-gam?
\choice
{$93{,}6\text{ kg}$}
{\True $9{,}36\text{ kg}$}
{$0{,}936\text{ kg}$}
{$936\text{ kg}$}
\loigiai{
Đổi: $1{,}2\text{ dm}^3 = 1\,200\text{ cm}^3$.\\
Khối lượng của khối kim loại là:
$1\,200 \times 7{,}8 = 9\,360\text{ (g)} = 9{,}36\text{ kg}$.\\
\textbf{Đáp án của câu hỏi là B}
}
\end{ex}

\begin{ex} %[5H8V46]
Một can dầu đang chứa $3\,500\text{ cm}^3$ dầu ăn. Người ta rót thêm vào can $1{,}8\text{ l}$ dầu ăn nữa thì vừa vặn đầy can. Dung tích chứa của can dầu đó là
\choice
{$4{,}3\text{ l}$}
{\True $5{,}3\text{ l}$}
{$5{,}8\text{ l}$}
{$6{,}3\text{ l}$}
\loigiai{
Đổi: $3\,500\text{ cm}^3 = 3{,}5\text{ dm}^3 = 3{,}5\text{ l}$.\\
Dung tích của can dầu là:
$3{,}5 + 1{,}8 = 5{,}3\text{ (l)}$.\\
\textbf{Đáp án của câu hỏi là B}
}
\end{ex}

\begin{ex} %[5H8V46]
Một bình nước đang chứa $18\text{ dm}^3$ nước, lượng nước này chiếm $75\%$ dung tích chứa của bình. Dung tích chứa của bình nước đó là
\choice
{$20\text{ dm}^3$}
{$22\text{ dm}^3$}
{\True $24\text{ dm}^3$}
{$25\text{ dm}^3$}
\loigiai{
Dung tích của bình nước là:
$18 : 75\% = 18 : 0{,}75 = 24\text{ (dm}^3\text{)}$.\\
\textbf{Đáp án của câu hỏi là C}
}
\end{ex}

\begin{ex} %[5H8V46]
Phép so sánh nào dưới đây là ĐÚNG?
\choice
{$0{,}8\text{ dm}^3 = 80\text{ cm}^3$}
{$1\,200\text{ cm}^3 > 1{,}5\text{ dm}^3$}
{\True $3\text{ dm}^3\text{ }5\text{ cm}^3 = 3{,}005\text{ dm}^3$}
{$4{,}2\text{ dm}^3 = 4\,020\text{ cm}^3$}
\loigiai{
Vì $1\text{ dm}^3 = 1\,000\text{ cm}^3$ nên $5\text{ cm}^3 = 0{,}005\text{ dm}^3$.\\
Do đó $3\text{ dm}^3\text{ }5\text{ cm}^3 = 3{,}005\text{ dm}^3$ là khẳng định đúng.\\
\textbf{Đáp án của câu hỏi là C}
}
\end{ex}

\begin{ex} %[5H8V46]
Một chai dung dịch có thể tích $0{,}25\text{ l}$. Mỗi lần sử dụng người ta dùng hết $25\text{ cm}^3$ dung dịch đó. Chai dung dịch đó dùng được trong bao nhiêu lần?
\choice
{$5\text{ lần}$}
{\True $10\text{ lần}$}
{$15\text{ lần}$}
{$20\text{ lần}$}
\loigiai{
Đổi: $0{,}25\text{ l} = 0{,}25\text{ dm}^3 = 250\text{ cm}^3$.\\
Số lần sử dụng là:
$250 : 25 = 10\text{ (lần)}$.\\
\textbf{Đáp án của câu hỏi là B}
}
\end{ex}

\begin{ex} %[5H8V46]
Một khối rubik hình lập phương có cạnh $6\text{ cm}$. Thể tích của khối rubik đó tính theo đơn vị đề-xi-mét khối là
\choice
{$21{,}6\text{ dm}^3$}
{$2{,}16\text{ dm}^3$}
{\True $0{,}216\text{ dm}^3$}
{$0{,}0216\text{ dm}^3$}
\loigiai{
Thể tích của khối rubik là:
$6 \times 6 \times 6 = 216\text{ (cm}^3\text{)}$.\\
Đổi: $216\text{ cm}^3 = 0{,}216\text{ dm}^3$.\\
\textbf{Đáp án của câu hỏi là C}
}
\end{ex}

\textbf{Phần 2. Trắc nghiệm đúng sai}
\setcounterpageref{ex}{0}

\begin{ex} %[5H8V46]
Cho ba số đo thể tích: $V_1 = 2{,}4\text{ dm}^3$; $V_2 = 2\,400\text{ cm}^3$; $V_3 = 2\text{ dm}^3\text{ }40\text{ cm}^3$.
\choiceTF
{\True Hai số đo $V_1$ và $V_2$ bằng nhau}
{\True Số đo $V_3$ đổi ra đề-xi-mét khối bằng $2{,}04\text{ dm}^3$}
{\True Thể tích $V_1$ lớn hơn thể tích $V_3$}
{Tổng thể tích $V_1 + V_3$ bằng $4{,}8\text{ dm}^3$}
\loigiai{
\begin{itemchoice}
\item Đúng. $V_1 = 2{,}4\text{ dm}^3 = 2\,400\text{ cm}^3 = V_2$.
\item Đúng. $V_3 = 2 + 40 : 1\,000 = 2{,}04\text{ dm}^3$.
\item Đúng. $2{,}4\text{ dm}^3 > 2{,}04\text{ dm}^3$.
\item Sai. $V_1 + V_3 = 2{,}4 + 2{,}04 = 4{,}44\text{ (dm}^3\text{)} 
eq 4{,}8\text{ dm}^3$.
\end{itemchoice}
}
\end{ex}

\begin{ex} %[5H8V46]
Một bể cá hình hộp chữ nhật có kích thước lòng trong là: chiều dài $50\text{ cm}$, chiều rộng $40\text{ cm}$ và chiều cao $30\text{ cm}$.
\choiceTF
{\True Thể tích của lòng bể cá là $60\,000\text{ cm}^3$}
{\True Thể tích đó tương đương với $60\text{ dm}^3$, bể chứa được tối đa $60\text{ l}$ nước}
{\True Khi mực nước trong bể cao $20\text{ cm}$, lượng nước có trong bể là $40\text{ l}$}
{Cần đổ thêm vào bể đúng $30\text{ l}$ nước nữa để đầy bể}
\loigiai{
\begin{itemchoice}
\item Đúng. Thể tích lòng bể: $50 \times 40 \times 30 = 60\,000\text{ (cm}^3\text{)}$.
\item Đúng. $60\,000\text{ cm}^3 = 60\text{ dm}^3 = 60\text{ l}$.
\item Đúng. Lượng nước hiện có: $50 \times 40 \times 20 = 40\,000\text{ (cm}^3\text{)} = 40\text{ dm}^3 = 40\text{ l}$.
\item Sai. Lượng nước cần đổ thêm là: $60 - 40 = 20\text{ (l)} 
eq 30\text{ l}$.
\end{itemchoice}
}
\end{ex}

\textbf{Phần 3. Trắc nghiệm trả lời ngắn}
\setcounterpageref{ex}{0}

\begin{ex} %[5H8V46]
Một hộp bánh hình lập phương có cạnh $10\text{ cm}$. Thể tích của hộp bánh đó bằng bao nhiêu đề-xi-mét khối?
\shortans{1}
\loigiai{
Thể tích hộp bánh là:
$10 \times 10 \times 10 = 1\,000\text{ (cm}^3\text{)}$.\\
Đổi: $1\,000\text{ cm}^3 = 1\text{ dm}^3$.
}
\end{ex}

\begin{ex} %[5H8V46]
Một thùng nước chứa $12\text{ l}$ nước. Người ta múc nước từ thùng ra các ca nhỏ, mỗi ca chứa $750\text{ cm}^3$. Hỏi múc được bao nhiêu ca nước đầy như thế?
\shortans{16}
\loigiai{
Đổi: $12\text{ l} = 12\text{ dm}^3 = 12\,000\text{ cm}^3$.\\
Số ca nước múc được là:
$12\,000 : 750 = 16\text{ (ca)}$.
}
\end{ex}

\begin{ex} %[5H8V46]
Một thanh kim loại hình hộp chữ nhật có kích thước $5\text{ cm} \times 8\text{ cm} \times 15\text{ cm}$. Thể tích của thanh kim loại đó là bao nhiêu xăng-ti-mét khối?
\shortans{600}
\loigiai{
Thể tích thanh kim loại là:
$5 \times 8 \times 15 = 600\text{ (cm}^3\text{)}$.
}
\end{ex}

\begin{ex} %[5H8V46]
Một bình nước hình hộp chữ nhật có đáy là hình vuông cạnh $10\text{ cm}$ và chiều cao $25\text{ cm}$. Người ta rót hết nước trong bình đầy vào các cốc loại $250\text{ cm}^3$. Hỏi rót được bao nhiêu cốc nước đầy?
\shortans{10}
\loigiai{
Thể tích của bình nước là:
$10 \times 10 \times 25 = 2\,500\text{ (cm}^3\text{)}$.\\
Số cốc nước rót được là:
$2\,500 : 250 = 10\text{ (cốc)}$.
}
\end{ex}

\begin{ex} %[5H8V46]
Người ta xếp $500$ khối lập phương nhỏ $1\text{ cm}^3$ khít nhau vào một chiếc hộp. Tổng thể tích của $500$ khối lập phương đó là bao nhiêu xăng-ti-mét khối?
\shortans{500}
\loigiai{
Tổng thể tích của $500$ khối lập phương là:
$500 \times 1 = 500\text{ (cm}^3\text{)}$.
}
\end{ex}

\begin{ex} %[5H8V46]
Một bể nước đang chứa $15\text{ l}$ nước. Người ta đổ thêm vào bể 3 ca nước, mỗi ca chứa đúng $1\,000\text{ cm}^3$. Lúc này trong bể có tất cả bao nhiêu lít nước?
\shortans{18}
\loigiai{
Thể tích 3 ca nước là:
$3 \times 1\,000 = 3\,000\text{ (cm}^3\text{)} = 3\text{ dm}^3 = 3\text{ l}$.\\
Tổng số lít nước trong bể là:
$15 + 3 = 18\text{ (l)}$.
}
\end{ex}
